\documentclass[a4paper,11pt,titlepage]{article}
\parskip 3pt

%% %%%%%%%%%%%%%%%%%%%% BEGIN PACKAGES %%%%%%%%%%%%%%%%%%%%

%%\usepackage[top=1in, bottom=1in, left=1in, right=1in]{geometry}

\usepackage{fullpage}
\usepackage{url}
\usepackage{hyperref}
\usepackage{listings}
\hypersetup{pdfborder=0 0 0}

%% Images
\usepackage{graphicx}

%% Side-by-side images
\usepackage{subfig}

%% Wrapped figures
\usepackage{wrapfig}

%% Drawing trees and other stuff
\usepackage{tikz}
\usetikzlibrary{trees,arrows}
\usepackage{amssymb}

%% Pseudocode
\usepackage[noend]{algorithmic}
\algsetup{indent=1.5em}

%% Linux Libertine
\usepackage[T1]{fontenc}
\usepackage{libertine}
\renewcommand*\oldstylenums[1]{{\fontfamily{fxlj}\selectfont #1}}

\let\stdhref\href
\renewcommand{\href}[2]{\stdhref{#1}{\texttt{#2}}}

\newcommand{\mailto}[1]{\href{mailto:#1}{#1}}

\let\stdsection\section         % because LaTeX cannot handle
                                % recursive commands
\renewcommand{\section}{\newpage\stdsection}

\let\tikzsquare\square
\renewcommand{\square}{\ensuremath\tikzsquare}

\newcommand{\code}[1]{\texttt{#1}}
\newcommand{\buzz}[1]{\emph{#1}}
\newcommand{\magic}[1]{\buzz{#1}}

\newcommand{\myhref}[1]{\href{http://#1}{#1}}

%% Clubs and widows. Avoid them like the plague.
\clubpenalty = 10000
\widowpenalty = 10000
\displaywidowpenalty = 10000

%% %%%%%%%%%%%%%%%%%%%% END COMMANDS %%%%%%%%%%%%%%%%%%%%

%Visigoth is a tool to generate, analyse and visualise Small World
%Networks. These are particular kinds of graphs which \emph{look} like naturally
%grown networks and share many mathematical properties with them.  For instance,
%friendship graphs and graphs of followers in modern web-based ``social
%networks'' are Small World Networks.

%As an educational tool Visigoth allows users to explore Small World Networks and
%hence can be used as an introduction to them. By exploring these networks users
%will have a firmer, more intuitive grasp of the mathematical properties they
%exhibit.

%The networks may be generated using several algorithms, allowing users to
%compare their output and its adherence to the properties of Small World
%Networks. The network generation algorithms in Visigoth form a smooth
%progression in the number of exhibited desired properties from completely random
%graphs, generated by the Erd\H{o}s-R\'{e}nyi algorithm, to fully compliant Small
%World Networks, generated by the Preferential Attachment
%algorithm. Additionally, Visigoth can connect to social networking platforms
%such as \Twitter to fetch and display real-world networks.

%We should keep it 2 pages long at most and should have the same structure as
%the Report


\begin{document}
\title{\includegraphics[width=2cm]{v.png}\\
  \Huge Visigoth\\\Large Graph visualisations}
\author{
  Andreea-Ingrid Funie\\\mailto{aif109@doc.ic.ac.uk}\and
  Alexandru Scvor\c tov\\\mailto{as10109@doc.ic.ac.uk}\and
  Francesco Mazzoli\\\mailto{fm2209@doc.ic.ac.uk}\and
  Marc-David Haubenstock\\\mailto{mh808@doc.ic.ac.uk}\and
  Maximilian Staudt\\\mailto{ms9109@doc.ic.ac.uk}
}
\date{January 2012}
\maketitle

\begin{description}
	\item[Introduction]
	% High Level, Nontechnical Description Why you should buy this product/listen to
	% this presentation? What is the functionality of the product?

    Visigoth is a tool to generate, analyse and visualise Small World
    Networks. These are particular kinds of graphs which \emph{look} like
    naturally grown networks and share many mathematical properties with them.
    For instance, friendship graphs and graphs of followers in modern web-based
    ``social networks'' are Small World Networks.

	\item[Visigoth's functionality]
	  %should we add here what is able to do?

    As an educational tool Visigoth allows users to explore Small World Networks
    in both 2D and 3D and hence can be used as an introduction to them. By
    exploring these networks users will have a firmer, more intuitive grasp of
    the mathematical properties they exhibit.

    The networks may be generated using several algorithms, allowing users to
    compare their output and its adherence to the properties of Small World
    Networks. The network generation algorithms in Visigoth form a smooth
    progression in the number of exhibited desired properties from completely
    random graphs, generated by the Erd\H{o}s-R\'{e}nyi algorithm, to fully
    compliant Small World Networks, generated by the Preferential Attachment
    algorithm. Additionally, Visigoth can connect to social networking platforms
    such as Twitter to fetch and display real-world networks.

    Users have easy access to common actions such as graph regeneration and node
    reshuffling. They can move the graph around, rotate it, generate node
    specific statistics and highlight direct node's neighbours.

    Users can also customize the appearance of the generated graph, such as:
    change edges and nodes colour, choose a different background, highlight
    cliques or highlight selected node and connections with it's direct
    neighbours.

    Finally, Visigoth allows taking snapshots of the graphs and exporting them
    to image files. The supported formats include PNG and JPEG, which are the
    image formats supported by most computers.

	\item[Technical Description]
	  %B. Short Technical Description; Short introduction into technologies used
	  %Design of your software, possibly including a diagram of the major
	  %components of the project; Main achievements

	\item[Software Engineering Issues]
	  %C. Software Engineering Issues: What technology was used and why; what
	  %other technology was considered but not used and why; Any technical
	  %challenges encountered and how addressed; Any risks anticipated, and how
	  %mitigated? Any collaboration/coordination difficulties encountered and how
	  %addressed; Development and testing methods and/or tools used; comparison of
	  %plans with actual achievements; Estimates of length of code in each of the
	  %components, or any other comparable measure of the effort required; Summary
	  %of each team member's contributions

    We encountered a number of difficulties whilst building Visigoth, which we
    over- came by dividing the application into mostly-independent components.
    The following sections provide overviews and short descriptions for each
    component.

    \begin{description}
      \item[Tests] The test system is heavily oriented towards functionality
        testing, checking that user actions cause the expected behaviour.

      \item[GUI] The GUI is the only user-visible component. It was relies
        heavily on Qt's widget library and OpenGL.

      \item[GraphScene] We took an Model-View-Controller approach in Visigoth's
        design. Whist the GUI system is the view, the GraphScene is the model
        and encap- sulates the controller.

      \item[Algorithms] The last component of Visigoth are the algorithms. All
        the individual algorithms implement the Algorithm interface, which
        allows them to be easily plugged into the GraphScene.
    \end{description}

    Work was divided between team members such that each would focus on
    developing a particular component and reviewing another. We chose an
    iterative development model, where each iteration consisted of adding a
    small number of features to the existing product.

    We have had to ensure that no bugs crept into the master branch. Towards
    this goal, no changes are merged into master until at least another member
    of the team reviews them. Through this practice of peer review the whole
    team has stayed informed on the current state of the project in addition to
    exposing and fixing bugs missed by the original code author.

    To further promote group awareness of the complete state of the project we
    held frequent short meetings true to the idea of \emph{scrum}. These were
    whole-group meetings held every few days, in which team members shared
    comments on what features they were working on at the time and the future
    outlook of the project.

    While our development methodology was not test-driven or focused on
    usability, we continuously tested the software on volunteers and towards the
    end of each iteration, we held meetings with our client. These tests and
    meetings uncovered a few problems, which eventually lead to small redesigns.

	\item[Validation]
	  %D. Validation and Conclusions How did you validate your product?

	During our product validation we distinguished between \emph{the client} (i.e.
	the person ``paying'' for the development of Visigoth) and \emph{the users}
	(i.e. the people who would end up using Visigoth).

	We held meetings with \emph{the client} near the end of each iteration
	in order to report on our progress and to get feedback on new
	developments. These meetings would consist of a short demo of newly
	added features, a short ``impromptu'' usability test and a discussion
	on what new features should be the focus of the following iteration.

	The feedback we received at the end of each iteration was consistently
	positive. The Visigoth interface always met our client's demands and
	only minor changes have been made from iteration to iteration

	Furthermore, we asked classmates and friends to try
	out accomplishing certain tasks in Visigoth while we watched. At the
	end, we would ask what felt awkward or wrong.

	The lasting impression was that they were happy with the software, but
	watching them struggle with certain aspects led to improvements in
	usability.

	\item[Conclusions]
	  %Was the project successful? What did you learn? What might you have done
	  % differently?

    Our team has thoroughly enjoyed the process that led to the creation of
    Visigoth. We managed to get the first results very quickly and then added
    features and refined the product in the successive months. As of today, we
    are very proud of the result and we plan to release it to the wider public
    in the near future.

    While we could have improved in some areas, the development of Visigoth has
    been very successful in what matters more: constantly delivering, on time,
    the features requested by the client while maintaining a clean and
    maintainable code base. We went even further than that and proposed features
    of our own which were all implemented. For this reason we are satisfied not
    only with the end result, but also with how smooth the development process
    was, and we think our client is satisfied as well.

\end{description}

\end{document}
